\section{Background} \label{sec:background}
The term prompt injection attack does not describe a specific attack against an LLM but rather a family or group of related attacks with similar modes of action and threat models.
In this section, we introduce some general definitions that can be applied to any prompt injection attack.
The definitions are mainly based on the formalization of prompt injection attacks by \liu
We will go into more details on specific prompt injections in the following section and discuss how the formalization of \liu relates to pre-existing prompt injection frameworks in~\cref{add_later}.

\subsection{LLM-Integrated Applications} \label{sec:llm_applications}
\begin{figure}[htb]
    \centering
    \includegraphics[width=\linewidth]{add_later.png}
    \caption{}
    \Description[<>]{<>}
    \label{fig:llm_application}
\end{figure}

An LLM-integrated application ties together three external components with the goal of fulfilling a target task.
The first component is an LLM $f$, which is controlled and managed by the application.
The second component is a user interacting with the application.
It is important to notice that, contrary to the colloquial term "using an LLM", the user does not directly access the LLM but only interacts with it through the LLM-integrated application.
The third component is any external data $\textbf{x}^t$ that is processed by the application.
This data is often, but not exclusively, text-based and can incorporate a wide range of information types like web pages, source code, and different document types.\\

The basic functionality of an LLM-integrated application is to utilize the LLM $f$ in order to process external data on behalf of the user.
To do so, the LLM is provided with a prompt $\textbf{p}$ consisting of textual instructions $\textbf{s}^t$ and external data $\textbf{x}^t$.
While the simplest way to combine instructions and data is through simple concatenation in the form of prefix prompts, there are other possible approaches.
In general, the data can be joined with the instructions through slot-filling with an arbitrary number and distribution of slots within the instructions.
We denote the arbitrary combination of instruction and data into a prompt with the notation shown in~\cref{eq:prompt_constr}.
\begin{equation}
    \textbf{p} = \textbf{s}^t \circ \textbf{x}^t
    \label{eq:prompt_constr}
\end{equation}
The instructions themselves can consist of parts originating from the user and the application.
Instructions from the application $\textbf{s}^a$, often referred to as the system prompt, usually define guard rails or steer the model towards a specific task that is intended by the application.
User instructions $\textbf{s}^u$ might contain a more specific task description, directly provided information, or in-context learning examples.
Analogously to the construction of the prompt, user and application instructions can be combined arbitrarily: $\textbf{s}^t = \textbf{s}^a \circ \textbf{s}^u$.
\\

The prompt $\textbf{p}$ is processed by the application's internal LLM and converted into the response $\textbf{r} = f(\textbf{p})$ that is returned to the user.
\Cref{fig:llm_application} depicts the connections between the application, LLM, user, and external data.
This high-level definition of an LLM-integrated application allows it to cover diverse real-world applications while still providing a common framework for prompt injection attacks.
Such real-world applications include the GUIs of popular LLMs like ChatGPT, Gemini, and Claude, as well as third-party services and even the incorporation of LLMs in application processes without the user being aware of their use.

\subsection{Prompt Injection Attacks}
Based on the formal definition of LLM-integrated applications in~\cref{sec:llm_applications}, prompt injection attacks can be defined.
The basic threat model of any kind of prompt injection attack includes an attacker pursuing the goal of changing the LLM's output from fulfilling the target task to fulfilling an attacker-defined injected task instead.
This is achieved by manipulating the external data that is processed by the application.
The capability of manipulating the external data is therefore the basic attacker capability that is present in every prompt injection threat model.
Intuitively, prompt injection attacks can be understood as an attacker exploiting the fact that the external data processed by LLM-integrated applications is often prone to manipulation.\\

Formally, the attacker alters the external target data $\textbf{x}^t$ to a compromised data set $\tilde{\textbf{x}}$.
He can do so not only by involving injected data $\textbf{x}^e$ but also injected instructions $\textbf{s}^e$.
Since there are diverse ways of altering the target data $\textbf{x}^t$, we denote the construction of the compromised data set through the execution of an attack function $\mathpzc{A}$ as shown in~\cref{eq:compr_data}.
\begin{equation}
    \tilde{\textbf{x}} = \mathpzc{A}(\textbf{x}^t, \textbf{x}^e, \textbf{s}^e)
    \label{eq:compr_data}
\end{equation}
In the case of an attack, the compromised external data $\tilde{\textbf{x}}$ is used to construct a compromised prompt $\tilde{\textbf{p}}$ analogously to~\cref{eq:prompt_constr}:
\begin{equation}
    \tilde{\textbf{p}} = \textbf{s}^t \circ \tilde{\textbf{x}}
    \label{eq:compr_prompt_constr}
\end{equation}
Considering the compromised data $\tilde{\textbf{x}}$ contains the injected data $\textbf{x}^e$ and more importantly, the injected instructions $\textbf{s}^e$, shows how the attacker takes influence on the LLM.
This demonstrates that prompt injection attacks are fundamentally based on the mixing of instructions and data in the construction of prompts.
We further discuss this fundamental problem in~\cref{add_later}\\

The general definition of prompt injection attacks by \liu only describes the fundamental characteristic shared amongst all prompt injection variants.
Two main parameters create the variety of prompt injection attacks found in the literature.
One parameter is the injected task defined by the attacker, the spectrum of which can be as broad as the target tasks of various applications.
A prominent example is the leaking of the system prompt by the LLM.
The other parameter is the definition of the attack function $\mathpzc{A}$.
Just like with the construction of prompts shown in~\cref{eq:prompt_constr}, the naive approach is to concatenate the instruction data with the components injected by the attacker.
However, much more complex approaches can be used, which can be adjusted to the attacked application's LLM, target task, and the injected task.
By presenting and discussing several case studies in the following chapters, we provide an overview of the variety of prompt injection techniques with different injected tasks and attack approaches.